%% kapa: exam=1 round=22 year=2011 period=2 form=A subject=accounting subject_ko=회계학 seq=24 number=24
%% source: Q-Net 감정평가사 기출문제 (cst00309) artlSeq=5036790 file=extracted/1008/5036790/2011년 제22회 감정평가사 1차시험 기출문제/2교시.hwp
%% conversion: kordoc 4.24.51 -> tools/gen1.py (auto); review=proofread-gpt-astra; flags=none
\begin{question}{24}
(주)서울은 20×1년 1월 1일에 영업사원 200명에게 각각 주식선택권 100개를 부여하였다. 각 주식선택권은 종업원이 앞으로 4년간 근무할 것을 조건으로 한다. 부여일 현재 주식선택권의 단위당 공정가치는 ￦5,000으로 추정되었다. (주)서울은 종업원 중 30\%가 4년 이내에 퇴사하여 주식선택권을 상실할 것으로 추정하였고, 실제로 20×1년과 20×2년에 각각 15명이 퇴사하였다. 동 주식선택권과 관련하여 (주)서울의 20×2년 포괄손익계산서상 당기손익에 반영될 주식보상비용은?
\begin{choices}
  \item ￦12,500,000
  \item ￦15,500,000
  \item ￦17,500,000
  \item ￦19,500,000
  \item ￦21,500,000
\end{choices}
\end{question}
\answer{3}
